% Part: lambda-calculus
% Chapter: lambda-definability
% Section: minimization

\documentclass[../../../include/open-logic-section]{subfiles}

\begin{document}

\olfileid[hi]{lam}{ldf}{min}
\olsection{न्यूनीकरण}

सामान्य पुनरावर्ती फलन वे हैं जिन्हें मूल फलनों $\Zero$, $\Succ$,
$\Proj{n}{i}$ से संयोजन, आदिम पुनरावर्तन और नियमित न्यूनीकरण द्वारा प्राप्त
किया जा सकता है। सभी सामान्य पुनरावर्ती फलनों को !!{lambda definable}
दिखाने के लिए हमें दिखाना होगा कि किसी !!a{lambda definable} फलन से नियमित
न्यूनीकरण द्वारा परिभाषित कोई भी फलन स्वयं !!{lambda definable} है।

\begin{lem}
  \ollabel{lem:min} यदि $f(x_1, \dots, x_k, y)$ नियमित और !!{lambda
  definable} है, तो निम्न द्वारा परिभाषित~$g$
  \[
  g(x_1, \dots, x_k) = \umin{y}{f(x_1,\dots,x_k, y) = 0}
  \]
  भी !!{lambda definable} है।
\end{lem}

\begin{proof}
  मान लें लैम्ब्डा पद~$F$ नियमित फलन $f(\vec x, y)$ को $\lambda$-परिभाषित
  करता है। $h$ को !!{lambda define} करने के लिए हम खोज-फलन और अचल-बिंदु
  संयोजक उपयोग करते हैं:
  \begin{align*}
    \fn{Search} & \ident \lambd[g][\lambd[f\,\vec{x}\,y][
        \fn{IsZero} (f\, \vec{x}\, y)\, y\, (g\, \vec{x} (\fn{Succ}\, y)]]\\
    H & \ident \lambd[\vec x][(Y \, \fn{Search}) F\, \vec{x}\, \num{0}],
  \end{align*}
  जहाँ $Y$ कोई भी अचल-बिंदु संयोजक है। अनौपचारिक रूप से $\fn{Search}$
  स्वसंदर्भी फलन है: ~$y$ से आरंभ करके जाँचता है कि $f\, \vec x\, y$
  शून्य है या नहीं; यदि हाँ तो~$y$ लौटाता है, अन्यथा $\fn{Succ}\, y$ के
  साथ स्वयं को पुकारता है। अतः $(Y \, \fn{Search}) F
  \num{n_1}\dots\num{n_k}\,\num{0}$ वह न्यूनतम~$m$ लौटाता है जिसके लिए
  $f(n_1, \dots, n_k, m) = 0$।
  
  विशेषतः ध्यान दें कि
  \begin{align*}
    (Y \, \fn{Search}) F \num{n_1}
    \dots\num{n_k}\, \num{m} & \red \num{m}
    \intertext{यदि $f(n_1, \dots,
      n_k, m) = 0$, अथवा}
    & \red (Y \, \fn{Search}) F\, \num{n_1} \dots\num{n_k}\, \num{m+1}
    \intertext{अन्यथा। चूँकि $f$ नियमित है, किसी $y$ के लिए $f(n_1,
      \dots, n_k, y) = 0$, और इसलिए}
    (Y \, \fn{Search}) F \num{n_1} \dots\num{n_k}\,\num{0}
    & \red \num{h(n_1, \dots, n_k)}.
    \end{align*}
\end{proof}


\begin{prop}
  प्रत्येक सामान्य पुनरावर्ती फलन !!{lambda definable} है।
\end{prop}

\begin{proof}
 \olref[prf]{lem:basic} से सभी मूल फलन !!{lambda definable} हैं; और
 \olref[prf]{lem:comp}, \olref[prf]{lem:prim} तथा \olref{lem:min} से
 !!{lambda definable} फलन संयोजन, आदिम पुनरावर्तन और नियमित न्यूनीकरण के
 अधीन संवृत हैं।
\end{proof}

\end{document}
